% id: 0496
% book: RPM 기하 · 05 벡터의 연산
% section: 유형 09 두 벡터가 서로 평행할 조건
% figure: none
% answer: $5$
% answer_source: 답지
% variant_level: 0 (원본 전사)
% 이 파일은 build.mjs 가 items.json 에서 만든다. 고칠 때는 items.json 을 고친다.
\smash{\raisebox{1.5mm}{\dmkichul{RPM 기하 0496번 · 79쪽 · 대표문제}}}
서로 평행하지 않고 영벡터가 아닌 두 벡터 $\vec{a}$, $\vec{b}$에 대하여 \par\smallskip\dispeq{$\vec{p}=\vec{a}+\vec{b}$, $\vec{q}=\vec{a}+3\vec{b}$, $\vec{r}=2\vec{a}+k\vec{b}$}\par\smallskip\noindent 일 때, 두 벡터 $\vec{p}+\vec{q}$, $\vec{r}-\vec{q}$가 서로 평행하도록 하는 실수 $k$의 값을 구하시오.
